\section{Uniform passages and dense visits}
\label{rc:section}

The covering argument needs a visit that creates unusually much sausage
area.  We construct such a visit by traversing a narrow serpentine ribbon.
The essential estimate is uniform over the history already drawn inside
the ribbon.  A positive probability for each fixed smooth template would
not give the required bound as the number of traversals increases.

Throughout this section $0<\kappa<8$ is fixed,
$d=1+\kappa/8$, and $p=2-d$.  For a curve segment $\eta$, write
$N_v(\eta)=\{z:\operatorname{dist}(z,\eta)<v\}$.
We use Corollary~\ref{geom:restart-support}, Proposition~\ref{ll:localization}, and Proposition~\ref{retain:sausage} in the following forms.  First, a smooth simple
slit can be followed in an open tube through ordered separating
crosscuts, with positive probability; the same assertion holds at a
stopped restart. Second, suppose a conformal map $F$ from $\mathbb H$ has a
continuous spherical extension to the closed half-plane and an access
curve from $0$ to modulus $b_0$ whose image is in $B(a,r/L)$, then
for all $L\ge L_0$,
\begin{equation}\label{rc:localization-input}
 \mathbb P\{F(\gamma[0,\sigma_{b_0L^b}])
                    \not\subset B(a,r)\}\le C L^{-c}.
\end{equation}
Here $\sigma_u$ is the first time the parameter-plane trace reaches modulus
$u$; the containment is interpreted up to infinite time if that stop is
not reached.  The constants $b,c,C>0$ and $L_0>1$ are uniform over the
maps and access curves, and the estimate includes parameter-boundary
visits. Third, for
each deterministic compact capacity-time interval $J$ of positive length,
\begin{equation}\label{rc:sausage-input}
 \lim_{a\downarrow0}\limsup_{v\downarrow0}
   \mathbb P\{v^{-p}\operatorname{Area}(N_v(\gamma(J)))<a\}=0.
\end{equation}
Only finite lists of deterministic intervals will be used.  In
particular, we do not condition this lower-tail estimate on a prescribed
passage or apply it directly to a random interval.

Rezaei's cover argument already uses conditional bounds after stopped
histories, compactness, and passages through prescribed spirals and tubes
\cite[Section~3, Lemmas~3.1--3.3]{rezaei2018}. Here we need one
probability bound for every past admitted by a ribbon step, together
with one deterministic small-mesh cutoff for those pasts. Iterating
this quantitative step will give the explicit lower bound
$c\exp(-Cm^2)$ for the dense visit.

\subsection{Geometry at the front of a ribbon}

Let $\sigma:[0,\infty)\to\overline{\mathbb H}$ be a unit-speed simple
curve, with Lipschitz tangent, curvature at most $1/10$, and normal $n$.
Assume that
\[
 P(s,x)=\sigma(s)+xn(s),\qquad s\ge0,\quad |x|\le1,
\]
is an embedded closed ribbon, lies in $\mathbb H$ when $s>0$, escapes
every bounded set as $s\to\infty$, and is straight vertical for
$0\le s\le4$, with $P(0,0)=0$.  The coordinate maps on patches of
length four have one fixed bi-Lipschitz constant. These geometric
bounds are part of the admissible class.  Fix $\epsilon=1/50$.

Consider any stopped chordal past that first reaches $s=k$, $k\ge1$,
has remained in $|x|<\epsilon$ up to that stop, meets the real boundary
only at the ribbon's base $P(\{0\}\times(-\epsilon,\epsilon))$, and ends at
$a=P(k,x_a)$ with $|x_a|\le\epsilon/4$.  The future portion of the
ribbon is in the surviving domain $D$: its interior is disjoint from the
past and has a route to the target along the ribbon.  The same holds for
the strips outside the core on each side of the preceding local patch.

\begin{lemma}[One ribbon step]\label{rc:step}
There are $q,c,v_0>0$, depending only on $\kappa$ and the fixed ribbon
constants, with the following property. For every stopped past
satisfying the preceding ribbon conditions and every deterministic
$0<v<v_0$, the conditional
probability is at least $q$ that the next first crossing of $s=k+1$
occurs in finite time, with tip offset $|x|<\epsilon/4$, while the
trace segment $\gamma_{\rm step}$ from the current
front to that next crossing stays in
\[
 k-1/4<s<k+1,\qquad |x|<\epsilon,
\]
before that crossing, and that
\begin{equation}\label{rc:step-mass}
 v^{-p}\operatorname{Area}
 \left(N_v(\gamma_{\rm step})\cap
   P((k+1/3,k+2/3)\times(-\epsilon,\epsilon))\right)\ge c.
\end{equation}
All three constants are uniform over the ribbon and the entire stopped
past.  For a ribbon of half-width $\ell$, the same assertion has mesh cutoff
$v_0\ell$ and lower bound $c\ell^d$ in \eqref{rc:step-mass}.
\end{lemma}

The goal is thus a single positive probability and a single mesh cutoff
for this entire class of pasts. We first establish uniform conformal
control of an access guide; the proof of Lemma~\ref{rc:step} then
transfers an interior sausage test through those conformal maps.

Use the guide
\[
 \alpha(t)=P(k+t,x_a(1-4t)_+),\qquad 0\le t\le2,
\]
and continue along the centerline to infinity.  It is a simple access
curve from the fresh tip to the target.  Freshness follows from the first
crossing of the front; the fresh-tip prime-end statement makes its
initial access agree with the driving prime end.  On this guide,
\begin{equation}\label{rc:guide-clearance}
 \operatorname{dist}(\alpha(t),\partial D)\ge c\min(t,1),
 \qquad |\alpha(t)-a|\le Ct\quad(0<t\le2).
\end{equation}
Indeed a local coordinate ball ahead of the front is free.  Its radius
can be chosen proportional to $t$ until the guide has advanced a fixed
distance, and fixed thereafter.  Nonlocal portions of the past cannot
enter these balls: they would meet the interior of this ribbon patch,
contrary to embedding.  Every ball is in the surviving component because
it connects to the future ribbon.

Set $w=P(k+1/2,0)$.  If $g:D\to\mathbb H$ sends the tip to $0$ and the
target to infinity, then
\begin{equation}\label{rc:waypoint-angle}
 \sin\arg g(w)\ge s_*>0.
\end{equation}
Here is a uniform harmonic-measure proof. The complete access guide
separates $D$ into two sides corresponding to the prime-end boundary
arcs mapped to the positive and negative real axes. An exit approached
from one side without crossing the guide uses that arc, even when its
physical exit point represents several prime ends. From $w$, on either
chosen side of the guide, take a route in the local ribbon to offset
$x=1/2$ or $x=-1/2$,
and backwards to a transverse rectangle centered at $s=k-1/2$.
These routes have uniformly bounded length and fixed positive clearance
from the core and the opposite side of the guide.  Across that rectangle
take a transverse route from the chosen offset to beyond the opposite
edge of the core.  The past has a longitudinal crossing of the rectangle:
take its portion between the last crossing of its rear edge preceding
the first crossing of its forward edge.  This connected crossing lies
in $|x|<\epsilon$ and separates the two transverse ends.  Thus a
continuous transverse crossing must meet the past before it reaches the
other end.

Starting at $w$, first run Brownian motion in a fixed small free ball
about $w$ until it reaches a smaller disk strictly on the chosen side of
the guide.  There is a uniformly positive probability to do so before
leaving the free ball.  Crossings of the guide during this initial move
are allowed.  From that smaller disk, Brownian motion has a uniformly
positive chance to follow the first route and then cross the transverse
rectangle.  One can specify this event in finitely many overlapping
coordinate disks and rectangular crossings with uniformly bounded aspect
ratios.  During this latter part, until the first domain exit, its route
avoids the access guide.  Its exit therefore belongs to the chosen
prime-end arc.  The strong Markov property gives a lower bound for the harmonic
measure of each arc from $w$.  In half-plane coordinates these two
harmonic measures are $\arg g(w)/\pi$ and $1-\arg g(w)/\pi$, proving
\eqref{rc:waypoint-angle}.  The argument uses no regularity of the past
inside the core.

After imposing $|g(w)|=1$, the conformal radius of $w$ is bounded above
and below, by \eqref{rc:guide-clearance} and the boundary point $a$.
Consequently the inverse maps have nondegenerate interior compact
limits.  The next two lemmas turn this interior normalization into a uniform
passage. First, the parameter guide approaches the driving point
uniformly. A small-image semicircle then joins an early localized trace
to the limiting guide. Finally, a finite choice of rational time
intervals transfers the sausage estimate. This will produce one
probability, mass threshold, and mesh cutoff valid for every past.
The estimates above verify the first lemma's hypotheses. Its limiting
map, guide and omitted boundary point then supply the hypotheses of
the second lemma.


\begin{figure}[tbp]
\centering
\begingroup
\definecolor{ribbonblue}{RGB}{31,88,138}
\definecolor{ribbonorange}{RGB}{181,85,22}
\begin{tikzpicture}[x=1cm,y=1cm,font=\small,
  line cap=round,line join=round,>=Latex]
\path[use as bounding box] (-.9,-1.05) rectangle (12.25,3.85);
\fill[ribbonblue!3] (0,0) rectangle (12,3);
\draw[black!40] (0,0)--(12,0) (0,3)--(12,3);
\fill[black!8] (0,1.15) rectangle (12,1.85);
\draw[black!35,dotted] (0,1.15)--(12,1.15) (0,1.85)--(12,1.85);
\fill[ribbonorange!20] (5.5,1.15) rectangle (7.5,1.85);
\draw[ribbonorange] (5.5,1.15) rectangle (7.5,1.85);
\draw[dashed,black!50] (3.5,-.2)--(3.5,3.3) (9.5,-.2)--(9.5,3.3);
\node[above] at (3.5,3.3) {front $s=k$};
\node[above] at (9.5,3.3) {$s=k+1$};
\draw[very thick] (0,1.4)--(.5,1.65)--(1,1.36)--(1.6,1.7)
 --(2.2,1.42)--(2.8,1.64)--(3.5,1.57);
\fill (3.5,1.57) circle (2pt);
\node[above left] at (3.5,1.7) {$a$};
\draw[thick,ribbonblue,->] (3.5,1.57)--(5,1.5)--(11.7,1.5);
\fill[white] (6.5,1.5) circle (2pt);
\draw[ribbonblue] (6.5,1.5) circle (2pt);
\node[below,align=center] at (6.5,.95) {waypoint\\$P(k+1/2,0)$};
\draw[thin,black!60] (6.5,.94)--(6.5,1.4);
\node[above,text=ribbonorange] at (6.5,2.25) {middle test region};
\draw[ribbonorange] (6.5,2.2)--(6.5,1.88);
\node[above,align=center] at (1.55,2.05) {past in the core\\$|x|<\epsilon$};
\node[above,text=ribbonblue] at (10.65,1.85) {access guide $\alpha$};
\node[left] at (0,3) {$x=1$};
\node[left] at (0,0) {$x=-1$};
\draw[<->] (3.5,-.65)--(9.5,-.65)
 node[midway,fill=white,inner sep=2pt] {one unit in $s$};
\end{tikzpicture}
\endgroup

\caption{\RibbonFrontCaption}
\label{fig:ribbon-front}
\end{figure}

Figure~\ref{fig:ribbon-front} shows the guide and the middle region
whose sausage mass appears in Lemma~\ref{rc:step}.


The compactness argument at the initial tip requires control of a guide,
but no common modulus of continuity for the inverse maps on the boundary.
We give that control explicitly.  Hyperbolic length in
$\mathbb H$ is normalized as $|dz|/\operatorname{Im}z$.

\Needspace{24\baselineskip}
\begin{lemma}[Compactness along a guide]\label{es:guide-compactness}
Let $F_n:\mathbb H\to D_n$ be conformal, with inverse $g_n$, and let
$\alpha_n:[0,T]\to\overline{D_n}$ be $C$-Lipschitz guides satisfying
\[
 \begin{gathered}
 \alpha_n(0)=a_n\in\partial D_n,\qquad
 \alpha_n(t)\in D_n\quad(t>0),\\
 \operatorname{dist}(\alpha_n(t),\partial D_n)
       \ge c\min(t,1).
 \end{gathered}
\]
Assume that $a_n$ stays in a bounded set and that
$\beta_n=g_n\circ\alpha_n$, initially defined for $t>0$, extends
continuously to $\beta_n(0)=0$.  Fix $t_w\in(0,T]$, put
$w_n=\alpha_n(t_w)$, and suppose that
\[
 |g_n(w_n)|=1,\qquad \operatorname{Im}g_n(w_n)\ge s>0.
\]
There is a subsequence on which $\alpha_n$ converges uniformly to a
guide $\alpha$, $F_n$ converges locally uniformly on $\mathbb H$ to a
nonconstant univalent map $F$, and $\beta_n$ converges locally uniformly
on $(0,T]$ to $\beta$, where $F\circ\beta=\alpha$.  In addition,
\begin{equation}\label{es:uniform-tip-collapse}
 \lim_{h\downarrow0}\sup_n\sup_{0\le t\le h}|\beta_n(t)|=0.
\end{equation}
Consequently $\beta(0+)=0$.  If $a_n\to a$, then $a\notin F(\mathbb H)$
and $F(\beta(t))\to a$ as $t\downarrow0$.
\end{lemma}

\begin{proof}
Write $z_n=g_n(w_n)$.  These points lie in a compact subset of
$\mathbb H$.  The conformal-radius identity and Koebe's theorem give
\[
 2\operatorname{Im}z_n\,|F_n'(z_n)|
   =R_{D_n}(w_n),\qquad
 c\min(t_w,1)\le R_{D_n}(w_n)\le4Ct_w.
\]
Here the upper bound uses $a_n\in\partial D_n$ and
$|w_n-a_n|\le Ct_w$.  Thus the inverse derivatives are bounded above
and away from zero at an interior compact set.  Since $w_n$ is bounded,
the growth and distortion theorems for normalized univalent functions
give a subsequence converging locally uniformly to a nonconstant
univalent $F$.  Arzel\`a--Ascoli also gives uniform convergence of the
physical guides.

For each fixed $\delta>0$, the part of $\alpha_n$ joining $w_n$ to any
$\alpha_n(t)$, $\delta\le t\le T$, has uniformly bounded hyperbolic
length.  Indeed the hyperbolic density in $D_n$ is at most twice the
reciprocal of the distance to its boundary, and the guides have speed
at most $C$.  Conformal invariance therefore confines
$\beta_n([\delta,T])$ to a fixed compact subset of $\mathbb H$.
The same estimates, or the conformal-radius identity for $g_n'$,
give equicontinuity there.  Subsequence limits satisfy
$F(\beta(t))=\alpha(t)$; univalence identifies them consistently for
all $\delta>0$.

It remains to prove \eqref{es:uniform-tip-collapse}.  Suppose otherwise.
There are $h_n\downarrow0$ and $\varepsilon>0$ such that
$\beta_n([0,h_n])$ reaches modulus $\varepsilon$.  After taking the
subsequence just described, there are a fixed disk $K\Subset\mathbb H$
and $r_0>0$ for which
\begin{equation}\label{es:waypoint-separation}
 |F_n(z)-a_n|\ge r_0\qquad(z\in K)
\end{equation}
for all sufficiently large $n$.  To see this, take $K$ sufficiently
small about the limit of $z_n$: the points $w_n$ have distance at least
$c\min(t_w,1)$ from $a_n$, and local convergence and distortion preserve
a smaller positive separation on $K$.

Put $T_n=\log(r_0/(Ch_n))$, which tends to infinity.  In the strip
$S=\mathbb R\times(0,\pi)$ set
\[
 U_n(v,\phi)=\min\left\{T_n,
       \log^+\frac{|F_n(e^{v+i\phi})-a_n|}{Ch_n}\right\}.
\]
Conformal invariance of Dirichlet energy and area change imply
\begin{equation}\label{es:collapse-energy}
 E_n:=\int_S|\nabla U_n|^2\,dv\,d\phi
 \le\int_{Ch_n<|w-a_n|<r_0}\frac{dA(w)}{|w-a_n|^2}
 =2\pi T_n.
\end{equation}
Every section $v<\log\varepsilon$ contains a zero of $U_n$:
the continuous curve $\beta_n([0,h_n])$ crosses that modulus, and
$|\alpha_n(t)-a_n|\le Ch_n$ throughout the corresponding physical
guide segment.  Angular Cauchy--Schwarz consequently gives, for almost
every such $v$,
\[
 \int_0^\pi U_n(v,\phi)^2\,d\phi
       \le\pi^2\int_0^\pi|\partial_\phi U_n(v,\phi)|^2\,d\phi.
\]
In particular this bounds the squared $L^2$ norm on the fixed unit
slab $(\log\varepsilon-1,\log\varepsilon)\times(0,\pi)$ by
$\pi^2E_n$.

For completeness, this bound propagates to every fixed bounded
rectangle $J\times(0,\pi)$ containing that slab.  Let
$m_n(v)=\pi^{-1}\int_0^\pi U_n(v,\phi)\,d\phi$.
There is a section $v_n$ in the initial slab with
$|m_n(v_n)|^2\le\pi E_n$, while
\[
 \int_J|m_n'(v)|^2\,dv\le E_n/\pi.
\]
Horizontal Cauchy--Schwarz bounds $|m_n(v)|^2$ on $J$ by a constant
depending only on $J$ and $\varepsilon$, times $E_n$.
The one-dimensional Poincar\'e inequality for
$U_n(v,\cdot)-m_n(v)$ then gives
\begin{equation}\label{es:anchored-ltwo}
 \int_{J\times(0,\pi)}U_n^2\le C_{J,\varepsilon}E_n.
\end{equation}
Choose $J$ so that $J\times(0,\pi)$ also contains $\log K$, using the
logarithm with argument in $(0,\pi)$.  By
\eqref{es:waypoint-separation}, $U_n=T_n$ on the
fixed positive-area set $\log K$.  Equations
\eqref{es:collapse-energy}--\eqref{es:anchored-ltwo} imply
\[
 |\log K|T_n^2\le C_{J,\varepsilon}2\pi T_n,
\]
a contradiction.  This proves uniform collapse.

Finally $F_n-a_n$ has no zeros in $\mathbb H$.  Hurwitz's theorem and
nonconstancy show that $F-a$ has no zeros there.  Since
$F(\beta(t))=\alpha(t)\to a$, the omitted point $a$ belongs to
$\partial F(\mathbb H)$.
\end{proof}

\begin{lemma}[A small image semicircle]\label{es:good-semicircle}
Let $F$ be univalent on $\mathbb H$, and let $a\notin F(\mathbb H)$.
Suppose $\beta:[0,T]\to\overline{\mathbb H}$ is continuous,
$\beta(0)=0$, $\beta(t)\in\mathbb H$ for $t>0$, and
$F(\beta(t))\to a$ as $t\downarrow0$.
Fix $r>0$ and $L>1$, and write $A=\log L$.  Choose $t_0>0$ so that
\[
 |F(\beta(t))-a|\le r/L\quad(0<t\le t_0),
 \qquad b_0=|\beta(t_0)|>0.
\]
In every unit logarithmic slab $(q,q+1)$ lying below $\log b_0$,
there is $B\in(e^q,e^{q+1})$ such that, provided $A>2\pi^2$,
\begin{equation}\label{es:semicircle-bound}
 \sup_{0<\theta<\pi}|F(Be^{i\theta})-a|
       \le r\exp\{-A+\pi\sqrt{2A}\}.
\end{equation}
Moreover
\begin{equation}\label{es:sector-bound}
 \sup_{\substack{B/8\le|z|\le B\\
                 \operatorname{Im}z\ge |z|/2}}
 |F(z)-a|
       \le r\exp\{-A+\pi\sqrt{2A}+6\}.
\end{equation}
For $A\ge64$ both upper bounds are smaller than $r/4$.
\end{lemma}

\begin{proof}
On $S=\mathbb R\times(0,\pi)$ use
\[
 U(v,\phi)=\min\left\{A,
          \log^+\frac{|F(e^{v+i\phi})-a|}{r/L}\right\}.
\]
As in \eqref{es:collapse-energy},
$\int_S|\nabla U|^2\le2\pi A$.  The initial guide supplies a zero on
every section $v<\log b_0$.  In the specified unit slab there is thus
a section $v=\log B$ whose angular energy is at most $2\pi A$.
Starting at a zero on that section and applying Cauchy--Schwarz gives
\[
 U(\log B,\phi)\le\sqrt{\pi\cdot2\pi A}
       =\pi\sqrt{2A}\qquad(0<\phi<\pi).
\]
The section may be chosen where the Sobolev restriction is absolutely
continuous; continuity in the open semicircle gives the bound at every
interior point.  Since $\pi\sqrt{2A}<A$, the upper truncation is
inactive and \eqref{es:semicircle-bound} follows.

Koebe's quarter theorem, applied after mapping the disk to
$\mathbb H$ at $z$, gives
\[
 \frac{\operatorname{Im}z}{2}|F'(z)|
 \le\operatorname{dist}(F(z),\partial F(\mathbb H))
 \le |F(z)-a|.
\]
Therefore $\log|F-a|$ is $2$-Lipschitz for hyperbolic distance.
If $z=\rho Be^{i\theta}$ is in the sector in
\eqref{es:sector-bound}, then $1/8\le\rho\le1$ and
$\sin\theta\ge1/2$, so
\[
 \cosh d_{\mathbb H}(z,iB)
   =\frac{1+\rho^2}{2\rho\sin\theta}
   \le\rho+\rho^{-1}\le65/8<\cosh3.
\]
Comparing with $iB$ in \eqref{es:semicircle-bound} proves
\eqref{es:sector-bound}.  For $A\ge64$,
$-A+\pi\sqrt{2A}+6<-\log4$.
\end{proof}

Here is the geometric consequence used to splice an initial passage
to the limiting guide.  Suppose $\beta$ in
Lemma~\ref{es:good-semicircle} is simple, and fix $\delta\in(0,T)$.
Choose the logarithmic slab so far to the left that its upper radius
is smaller than
\[
 \min\left\{b_0/2,\ \min_{\delta\le t\le T}|\beta(t)|\right\}.
\]
The last time $t_B\le T$ with $|\beta(t_B)|=B$ then satisfies
$t_B<\delta$, and the remaining guide lies strictly outside that
semicircle.  Put $b=B/4$.  If an initial hull stays in
\[
 \{z:|\operatorname{Re}z|<b/4,\quad
           0\le\operatorname{Im}z\le b\}
\]
and ends at $\zeta=x+ib$, its modulus is less than $B/2$.
From $\zeta$ move vertically upwards to the semicircle $|z|=B$,
follow the semicircle to $\beta(t_B)$, and then follow
$\beta([t_B,T])$.  The vertical piece belongs to the sector in
\eqref{es:sector-bound}; the semicircular piece obeys
\eqref{es:semicircle-bound}.  Both therefore map into $B(a,r/4)$
when $\log L\ge64$.  The vertical piece is above the initial hull,
and the other pieces are outside its containing disk.  These pieces
form a simple continuation route; corners can be smoothed in its
open neighborhoods.  Each fixed such route after $\zeta$ lies in a
compact subset of $\mathbb H$.  Consequently local uniform
convergence of the maps is sufficient to transfer all strict
continuation restrictions along that route.


\subsection{The uniform passage estimate}

We now prove Lemma~\ref{rc:step}. If uniform constants failed,
a sequence of admissible states and vanishing meshes would have a
nondegenerate conformal limit. A passage in that limit has positive
probability. Selecting finitely many compact sets, margins and
rational time intervals transfers it to the sequence; subtracting
the unconditional sausage exception preserves a fixed positive
passage-and-mass probability.




\begin{proof}[Proof of Lemma~\ref{rc:step}]
Suppose the uniform assertion were false.  Then for each integer $n$
there would be admissible data and $0<v_n<1/n$ such that the probability
of passage together with normalized mass at least $1/n$ is less than
$1/n$.  Denote their ribbon maps and front positions by $P_n$ and
$k_n$. Translate the physical starting tips to a bounded region, keeping
the same notation, and put
\[
 V=[-1/2,3/2]\times[-1,1],\qquad Q_n(u,x)=P_n(k_n+u,x).
\]
On a subsequence these charts converge uniformly, together with the
centerline tangents, to a chart $Q$ with the same bi-Lipschitz bound.
The limiting guide through relative coordinate $u=5/4$ is therefore
simple. Lemma~\ref{es:guide-compactness}, applied to this part of the
guide, gives locally uniform convergence of the normalized inverse maps
$F_n$ to a nonconstant univalent $F$. The parameter guides converge away
from time zero to a simple guide $\beta$ with $\beta(0+)=0$. Write
$a_n\to a$ for the physical tips.

Use one ordinary parameter-plane SLE $\gamma$ for all $n$. For each
fixed $n$, the image $F_n\circ\gamma$, after the corresponding capacity
rescaling, has the law of the conditional continuation in that state.
This realizes the events below on one probability space; no joint
relation between the original stopped states is needed.

Choose $r>0$ so small that, for all sufficiently large $n$ and also for
the limiting chart,
\[
 \begin{aligned}
 \overline B(a_n,2r)&\subset
 Q_n\bigl((-1/8,1/4)\times(-\epsilon/2,\epsilon/2)\bigr),\\
 \overline B(a,r)&\subset
 Q\bigl((-1/8,1/4)\times(-\epsilon/2,\epsilon/2)\bigr).
 \end{aligned}
\]
The centered tips and the common bi-Lipschitz bound permit this choice.
In particular, a trace in the first ball has not yet reached the middle
region or the next front.

We fix the probability of the initial vertical passage before choosing
the localization scale. Let $\tau_b$ be the first time ordinary
parameter-plane SLE attains height $b$. Slit support gives a constant
$A<\infty$ for which
\[
 E_b=\left\{\tau_b\le Ab^2,\quad
       |\operatorname{Re}\gamma(t)|<b/4\ (0\le t\le\tau_b)\right\}
\]
has probability $q_0>0$. Brownian scaling makes $q_0$ independent of
$b>0$, and $E_b\in\mathcal F_{\tau_b}$. Fix $L\ge L_0$ large enough that
$\log L\ge64$ and $CL^{-c}<q_0/2$, using the constants in
\eqref{rc:localization-input}.

A sufficiently short fixed physical guide segment lies in
$B(a_n,r/L)$ for every large $n$. Its endpoint has parameter modulus
bounded below by some $b_0>0$: this follows from interior convergence
at that fixed positive guide time. Decrease $b_0$ so that the same
statement holds for the limiting guide, and truncate the guides at
their first modulus-$b_0$ points when using them as localization
access curves. Keep the full limiting guide for the continuation.
Lemma~\ref{es:good-semicircle}
supplies a radius $B<b_0/2$ whose semicircle and central sector map
under $F$ into $B(a,r/4)$. Choose its logarithmic slab sufficiently
far to the left that the last crossing of $|z|=B$ by $\beta$ occurs
while its physical image is still in $B(a,r/4)$. The rest of the guide,
through $u=5/4$, then lies outside that semicircle.

Put $b=B/4$, $\tau=\tau_b$, and $E=E_b$. On $E$ the trace through
$\tau$ stays in
\[
 |\operatorname{Re}z|<b/4,\qquad 0\le\operatorname{Im}z\le b.
\]
Its filled hull stays there too, since the part of $\mathbb H$ outside
the closed rectangle is connected to infinity. The terminal tip is fresh
and has an upward vertical access ray. Let $G_n$ be the event that the
physical image under $F_n$ of this initial trace stays in $B(a_n,r)$.
It is $\mathcal F_\tau$-measurable. Lemma~\ref{geom:extension} gives
continuous boundary extensions of the inverse maps, so the short access
curve and \eqref{rc:localization-input} give
\[
 \mathbb P(E\setminus G_n)\le CL^{-c}.
\]
Indeed, on $E$ the trace through $\tau$ has modulus less than
$B/2<b_0$, so the localization estimate controls this entire prefix.
Our choice of $L$ yields $\mathbb P(E\cap G_n)\ge q_0/2$.

Given a realization in $E$, use the geometric splice after
Lemma~\ref{es:good-semicircle}: move upwards from its tip to the
semicircle of radius $B$, follow the semicircle to the last crossing
of $\beta$, and then follow $\beta$ through $u=5/4$. This simple route
avoids the previous hull. Its initial pieces map into the early ball
$B(a,r/4)$; thereafter it follows the guide, whose longitudinal
coordinate increases and whose offset is zero by $u=1/4$.

We describe a passage event with margins that persist under uniform
changes of the maps. For $\delta=1/\nu<\min(\epsilon/8,1/12)$, where
$\nu$ is a positive integer, define the compact coordinate sets
\[
 V_\delta=\left\{(u,x)\in V:
 \begin{gathered}
 -1/4+\delta\le u\le3/2-\delta,\quad |x|\le\epsilon-\delta,\\
 u\le1-\delta\ \text{or}\ |x|\le\epsilon/4-\delta
 \end{gathered}\right\},
\]
\[
 \begin{aligned}
 P_\delta&=V_\delta\cap\{u\le1-2\delta\},&
 O_\delta&=V_\delta\cap\{u\ge1+\delta\},\\
 M_\delta&=[1/3+\delta,2/3-\delta]
                  \times[-\epsilon+\delta,\epsilon-\delta].
 \end{aligned}
\]
The set $V_\delta$ keeps the whole suffix in the core and centers it
near the new front. The sets $P_\delta$, $M_\delta$, and $O_\delta$
specify a prefix before that front, a middle interval, and a point
beyond the front. The closed alternative in $V_\delta$ is useful:
these restrictions are closed conditions on a continuous path.

Let
\[
 K_j=\{z:|\operatorname{Re}z|\le j,\quad
                1/j\le\operatorname{Im}z\le j\},\qquad j=1,2,\ldots.
\]
Let $C$ be the union, over $j,\nu$ as above and rational
$0<t_-<t_+<t_{\rm out}$, of the following events:
\[
 \begin{gathered}
 \tau<t_-,\qquad \gamma([\tau,t_{\rm out}])\subset K_j,\\
 F(\gamma([\tau,t_{\rm out}]))\subset Q(V_\delta),\qquad
 F(\gamma([\tau,t_+]))\subset Q(P_\delta),\\
 F(\gamma([t_-,t_+]))\subset Q(M_\delta),\qquad
 F(\gamma(t_{\rm out}))\in Q(O_\delta).
 \end{gathered}
\]
A witness is false unless its $K_j$ containment holds, so $F$ is
evaluated only in $\mathbb H$. All target sets in the display are
compact. By continuity, the interval containments can be checked on
dense rational times, together with $\gamma(\tau)$. Each witness is
therefore $\mathcal F_{t_{\rm out}}$-measurable. Their countable union
$C$ depends on the fixed limit map and chart, independently of $n$,
$G_n$, and the route used to establish its probability.

For each realization in $E$, choose a sufficiently thin tube around
the spliced route. Its image under $F$ lies in the interior of
$Q(V_\delta)$ for some allowed $\delta$. Choose two ordered middle
gates so that the $F$-image of the tube portion from its initial tip
through the second gate lies in the interior of $Q(P_\delta)$, and
the $F$-image of the portion between the gates lies in the interior
of $Q(M_\delta)$. Choose the terminal parameter open set with
$F$-image in the interior of $Q(O_\delta)$. The parameter-plane closure
of the post-$\tau$ tube, including its initial tip, lies in some $K_j$.
Before its first passage through the second gate, a path following the
embedded tube remains in its prefix portion. It visits an open region
between the gates, so continuity gives a nondegenerate rational interval
$[t_-,t_+]$ there before that second passage. A later visit to the
terminal open set supplies a rational $t_{\rm out}>t_+$. These times
satisfy the displayed conditions. Map out the initial hull. Prime-end
access and continuity at the driving point permit a smooth simple slit in the
restarted half-plane following the tube. Slit support and the domain
Markov property give
\[
 q(\omega):=\mathbb P(C\mid\mathcal F_\tau)(\omega)>0
 \quad\text{on }E
\]
outside the usual null set. This argument needs no measurable choice
of a route, since $C$ itself is the fixed countable union just defined.
The value of $q$ may depend on the past.

Choose $q_1>0$ so that
$\mathbb P(E\cap\{q<q_1\})<q_0/4$. Since $G_n$ is
$\mathcal F_\tau$-measurable, for every large $n$,
\begin{equation}\label{rc:varying-good-events}
 \mathbb P(E\cap G_n\cap C)
 =\mathbb E[\mathbf1_{E\cap G_n}q]
 \ge q_1q_0/4=:q_2>0.
\end{equation}
Choose a finite subfamily of witnesses whose omission from $C$ has
probability less than $q_2/4$ on $E\cap C$. The same absolute loss
bound holds after intersecting with any of the varying events $G_n$.

We next transfer a retained witness to the $n$th chart. If $z$ is on
its suffix, write $F(z)=Q(v)$ with $v\in V_\delta$. Such $v$ has
distance at least $\delta$ from $\partial V$. If $K$ is the common
bi-Lipschitz constant, then $Q_n(v)$ has distance at least $\delta/K$
from the Jordan boundary $Q_n(\partial V)$. On the finitely many
retained sets,
\[
 \sup_{z\in K_j}|F_n(z)-F(z)|
       +\sup_{v\in V}|Q_n(v)-Q(v)|\longrightarrow0.
\]
For all sufficiently large $n$, the displayed sum is less than
$\delta/(4K)$ for every retained witness. Thus $F_n(z)$ lies in the
ball about $Q_n(v)$ of radius $\delta/K$. That ball stays inside
the Jordan domain $Q_n(\operatorname{int}V)$, so the inverse coordinates
are defined. The inverse Lipschitz bound then gives
\[
 \bigl|Q_n^{-1}(F_n(z))-v\bigr|<\delta/4
\]
uniformly along every retained suffix.

The prefix on $G_n$ has relative coordinate $u_n<1/4$, whereas the
point at $t_{\rm out}$ has $u_n>1$. Continuity forces a first crossing
of $u_n=1$ before $t_{\rm out}$. The full $V_\delta$ containment gives
the rear and core bounds until this crossing. At the crossing, the
limit coordinate satisfies $u\ge1-\delta/4>1-\delta$, so the closed
alternative in $V_\delta$ forces $|x_n|<\epsilon/4$. The
$P_\delta$ containment gives $u_n\le1-7\delta/4<1$ through $t_+$.
Thus $J=[t_-,t_+]$ precedes the first crossing. Its image lies inside
the $n$th middle region with a positive physical margin, by the
$M_\delta$ bounds and the common bi-Lipschitz control.

For each retained $K_j$, use the larger rectangle $K_{2j}$.
Derivative convergence gives deterministic constants
$0<c_*<C_*<\infty$ such that
\[
 c_*\le |F_n'|\le C_*
\]
there for all sufficiently large $n$. Fix $\lambda>0$ sufficiently
small. On a retained witness, since $v_n\to0$, the
$\lambda v_n$-sausage of its interval lies in $K_{2j}$; the segment
from a sausage point to its trace point also lies there. The upper
derivative bound maps this
sausage into the physical $v_n$-sausage, and the middle-region margin
keeps its image inside that region. Change of area gives
\begin{equation}\label{rc:compact-sausage-transfer}
 v_n^{-p}\operatorname{Area}(\hbox{physical sausage in the middle})
 \ge c_*^2\lambda^p
   (\lambda v_n)^{-p}\operatorname{Area}
      (N_{\lambda v_n}(\gamma(J))).
\end{equation}

Apply \eqref{rc:sausage-input} to the finite list of deterministic
intervals. It supplies a fixed positive lower mass threshold $a_*$ with
total failure probability less than $q_2/4$ at all sufficiently small
$v_n$. This is an unconditional estimate, so it can be subtracted
from \eqref{rc:varying-good-events}. After the finite extraction and
this sausage exception, the remaining probability is at least
$q_2-q_2/4-q_2/4=q_2/2$. On that event,
\eqref{rc:compact-sausage-transfer} gives the fixed mass threshold
$c_*^2\lambda^p a_*$. This event is contained in the step event stated
in the lemma: its sausage comes from $J$, which is before the first
crossing, and the transferred path has all the required bounds.
That step event is measurable from the trace stopped at the first
crossing. Although the auxiliary witness $C$ may inspect a later suffix
through $t_{\rm out}$, it is used only to prove this inclusion and
lower-bound the probability of the stopped event. The later suffix is
discarded before any iteration. For all sufficiently large $n$, both
the fixed mass threshold and $q_2/2$ exceed $1/n$, contradicting the
choice of the data.

This proves the existence of deterministic $q,c,v_0>0$ uniformly over
admissible states, at each prescribed deterministic mesh below $v_0$.
Scaling spatial coordinates by $\ell$ multiplies normalized sausage area by
$\ell^{2-p}=\ell^d$, proving the final assertion.
\end{proof}

\subsection{A dense visit in bounded conformal-radius time}

\begin{proposition}[Dense visit]\label{rc:dense}
Fix $s_0>0$. There are constants $c,C,C_2>0$, $C_1>1$, and an integer
$m_0$ with the following property. Let $S$ be a stopping time of ordinary
chordal SLE in $\mathbb H$, and let $T_z$ be the lifetime of an interior
marked point $z$. The assertion below holds conditionally almost surely
on $\{S<\infty,\ S<T_z\}$, in the surviving state at $S$ or its image
under a translation and positive dilation, with the transformed marked
point still denoted by $z$.
Suppose the marked point in that state has conformal radius $r$ and angle
sine at least $s_0$. Let $\rho$ be the first subsequent live stop at which
its conformal radius is $r/C_1$, with $\rho=\infty$ if that level is not
reached, and use the finite centered law of Lemma~\ref{ll:finite-tilt}
through $\rho$. For each integer $m\ge m_0$ there is a deterministic
$\delta_m>0$, independent of the stopped state, such that, for every
prescribed deterministic physical mesh $0<v\le r\delta_m$, there is an
$\mathcal F_\rho$-measurable success event with conditional centered
probability at least
\begin{equation}\label{rc:dense-probability}
 c\exp(-Cm^2)
\end{equation}
on which the experiment ends at a finite stopping time $T<\rho$ and
generates sausage mass at least
\begin{equation}\label{rc:dense-mass}
 v^{-p}\operatorname{Area}
  \bigl(N_v(\gamma_{\rm experiment})\cap B(z,C_2r)\bigr)
       \ge c r^d m^p.
\end{equation}
\end{proposition}

\begin{proof}
Normalize the stopped half-plane map so that $w=g(z)$ lies on the upper
unit semicircle.  Write $y=\operatorname{Im}w\ge s_0$.  Initially work
under the ordinary chordal law in these coordinates.  Let $w'$ be
$w$ displaced horizontally by $s_0/4$ on the side away from the imaginary
axis, choosing either side if $\operatorname{Re}w=0$.  The closed disk
of radius $s_0/20$ about $w'$ is inside $B(w,y/2)$ and stays a fixed
positive distance from $w$. Denote this disk by $\mathcal D$.

Set $h_0=s_0/200$.  In this disk place $2m+1$ vertical segments from
height $y-h_0$ to $y+h_0$, with equally spaced horizontal coordinates
in an interval of length $2h_0$ centered at $\operatorname{Re}w'$.
Join consecutive ends by exterior semicircles, alternately above and
below the array, and traverse the segments in serpentine order.  The
semicircles have radii comparable to $h_0/m$.  Approach the first segment
from below along a smooth simple connector from $0$, initially vertical,
whose horizontal displacement is made between heights $s_0/8$ and
$s_0/2$.  Continue the last segment vertically upwards to infinity.
The connector may be chosen from a family with uniform smooth bounds,
with $w$ a fixed positive distance from it; the sign choice gives two
such compact families as $w$ varies.

Choose
\[
 \ell=\frac{h_0}{C_0m}
\]
with $C_0$ a sufficiently large fixed constant.  The normal ribbon of
width $\ell$ about this entire guide is embedded and, in units of
$\ell$, has the geometry required above.  Indeed the nonadjacent
straight passes have separation comparable to $h_0/m\gg\ell$;
the turn radii are at least $10\ell$; and the smooth connector has a
fixed-width embedded ribbon on its original scale and straight
attachments.  The last segment and approach lie in different columns.
The ribbon starts with a vertical segment of length at least $4\ell$.
Let $s_{\rm out}$ be arclength in the normalized half-plane from the
base to the point where the final vertical continuation exits $\mathcal D$,
and set
\[
 N=\left\lceil s_{\rm out}/\ell\right\rceil+1.
\]
We end the experiment at the first crossing of the front at arclength
$N\ell$. This front lies between $\ell$ and $2\ell$ beyond the disk
exit. Since $s_{\rm out}=O(m)$ and $\ell\asymp m^{-1}$, the number of
unit steps satisfies $N\le Cm^2$. At least $cm^2$ middle strips before
this front lie wholly inside $\mathcal D$.

Figure~\ref{fig:serpentine-ribbon} shows the displaced disk and the
serpentine array within it.
\begin{figure}[!htbp]
\centering
\begingroup
\definecolor{ribbonblue}{RGB}{31,88,138}
\definecolor{ribbonorange}{RGB}{181,85,22}
\begin{tikzpicture}[x=1cm,y=1cm,font=\small,
  line cap=round,line join=round,>=Latex]
\path[use as bounding box] (-.2,-.55) rectangle (14.7,6.05);
\begin{scope}[shift={(0,-6)}]
\node[anchor=west,font=\small\bfseries] at (0,11.85)
  {(a) Marked point and displaced disk};
\node[anchor=west,font=\small\bfseries] at (7.25,11.85)
  {(b) The array inside the disk (enlarged)};

\draw[black!70] (.1,6.55)--(6.6,6.55);
\node[anchor=north east] at (6.55,6.49) {$\mathbb R$};
\node[anchor=west] at (.16,10.95) {$\mathbb H$};
\fill (1.65,9.65) circle (2pt);
\node[anchor=south east,align=right] at (1.56,9.7)
  {$w=g(z)$\\[-2pt]$|w|=1$};
\draw[black!45,dashed] (3.6,9.65) circle (1.03);
\fill[ribbonblue!4] (3.6,9.65) circle (1.03);
\fill[black!60] (3.6,9.65) circle (1.3pt);
\node[fill=white,inner sep=1.5pt] at (5.45,9.46) {$w'$};
\draw[thin,black!60] (5.22,9.49)--(3.65,9.63);
\draw[<->,thin] (1.68,10.92)--(3.6,10.92)
  node[midway,fill=white,inner sep=1.4pt] {$s_0/4$};
\draw[thin,black!45] (1.65,10.16)--(1.65,11.08);
\draw[thin,black!45] (3.6,10.68)--(3.6,11.08);
\node[anchor=west,align=left,font=\footnotesize] at (4.58,10.27)
  {radius\\$s_0/20$};
\draw[thin,black!50] (4.57,10.29)--(4.3,10.41);
\draw[ribbonblue,line width=1.1pt]
  (.65,6.55)--(.65,7.15)
  .. controls (.65,7.95) and (2.89,7.63) .. (2.89,8.5)
  --(2.89,9.08)
  --(2.89,10.05) arc[start angle=180,end angle=0,radius=.12]
  --(3.13,9.08) arc[start angle=180,end angle=360,radius=.12]
  --(3.37,10.05) arc[start angle=180,end angle=0,radius=.12]
  --(3.61,9.08) arc[start angle=180,end angle=360,radius=.12]
  --(3.85,10.05) arc[start angle=180,end angle=0,radius=.12]
  --(4.09,9.08) arc[start angle=180,end angle=360,radius=.12]
  --(4.33,10.05)--(4.33,11.27);
\draw[ribbonblue,-{Latex[length=2mm]}] (4.33,11.02)--(4.33,11.33);
\node[anchor=west,font=\footnotesize] at (4.43,11.18) {future guide};
\fill (.65,6.55) circle (1.8pt);
\node[anchor=north] at (.65,6.43) {$0$};
\node[anchor=north,align=center,font=\footnotesize] at (2.65,7.47)
  {smooth connector\\initially vertical};

\begin{scope}[shift={(7.45,6.92)}]
  \fill[ribbonblue!3] (-.15,-.08) rectangle (4.95,3.81);
  \node[anchor=west,font=\footnotesize,text=black!65] at (.03,3.7)
    {$2m+1$ vertical strands};
  \draw[ribbonblue!15,line width=4pt]
    (.2,-.32)--(.2,2.65) arc[start angle=180,end angle=0,radius=.26]
    --(.72,.5) arc[start angle=180,end angle=360,radius=.26]
    --(1.24,2.65) arc[start angle=180,end angle=0,radius=.26]
    --(1.76,.5) arc[start angle=180,end angle=360,radius=.26]
    --(2.28,2.65) arc[start angle=180,end angle=0,radius=.26]
    --(2.8,.5) arc[start angle=180,end angle=360,radius=.26]
    --(3.32,2.65) arc[start angle=180,end angle=0,radius=.26]
    --(3.84,.5) arc[start angle=180,end angle=360,radius=.26]
    --(4.36,3.37);
  \draw[ribbonblue,line width=.85pt]
    (.2,-.32)--(.2,2.65) arc[start angle=180,end angle=0,radius=.26]
    --(.72,.5) arc[start angle=180,end angle=360,radius=.26]
    --(1.24,2.65) arc[start angle=180,end angle=0,radius=.26]
    --(1.76,.5) arc[start angle=180,end angle=360,radius=.26]
    --(2.28,2.65) arc[start angle=180,end angle=0,radius=.26]
    --(2.8,.5) arc[start angle=180,end angle=360,radius=.26]
    --(3.32,2.65) arc[start angle=180,end angle=0,radius=.26]
    --(3.84,.5) arc[start angle=180,end angle=360,radius=.26]
    --(4.36,3.37);
  \foreach \x in {.2,.72,1.24,1.76,2.28,2.8,3.32,3.84,4.36}
    \foreach \y in {.87,1.23,1.59,1.95,2.31}
      \fill[ribbonorange] (\x-.037,\y-.043) rectangle (\x+.037,\y+.043);
  \draw[ribbonblue,-{Latex[length=1.7mm]}] (.2,.06)--(.2,.35);
  \draw[ribbonblue,-{Latex[length=1.7mm]}] (4.36,3.06)--(4.36,3.45);
  \draw[<->,thin] (.2,-.53)--(4.36,-.53)
    node[midway,fill=white,inner sep=1.4pt] {$2h_0$};
  \draw[<->,thin] (5.07,.5)--(5.07,2.65)
    node[midway,fill=white,inner sep=1.5pt,rotate=90] {$2h_0$};
  \node[anchor=west,text=ribbonorange,align=left,font=\footnotesize]
    at (5.38,2.85) {disjoint\\middle\\regions};
  \draw[ribbonorange,thin] (5.25,2.85)--(4.4,2.31);
\end{scope}
\node[anchor=north,font=\small,align=center] at (10.25,5.97)
  {$\ell=h_0/(C_0m)\asymp m^{-1},\qquad N\asymp m^2$};

\end{scope}
\end{tikzpicture}
\endgroup

\caption{\SerpentineRibbonCaption}
\label{fig:serpentine-ribbon}
\end{figure}

Let $F$ be the inverse of the original stopped map. Choose a fixed
$\chi>0$, depending only on $s_0$, so that the closed
$\chi$-neighborhood of $\mathcal D$ is inside $B(w,y/2)$. Such a
choice is uniform because the disk center is at distance $s_0/4$
from $w$ and its radius is $s_0/20$, while $y\ge s_0$.
On this neighborhood, Koebe distortion gives
\[
 c r\le |F'|\le C r,
 \qquad F(\hbox{that neighborhood})\subset B(z,C_2r).
\]
Here $|F'(w)|=r/(2y)$ and $s_0\le y\le1$; all constants are uniform
in the stopped domain. For the prescribed physical mesh $v$, put
\[
                     \eta=\lambda v/r,
\]
where $\lambda>0$ is small enough for the upper Lipschitz comparison.
Choose the deterministic $\delta_m>0$ so that
\[
 \lambda\delta_m<\min\{v_0\ell,\chi/2\}.
\]
Then $v\le r\delta_m$ makes the same parameter mesh $\eta$ admissible
at every ribbon step. It also keeps the segment from any sausage point
counted in $\mathcal D$ to a witnessing trace point inside the
distortion neighborhood.

Use vertical-slit support for the initial step from the real boundary,
requiring confinement and centered arrival only.  It has a fixed
positive probability by scaling.  At each subsequent first front
crossing apply Lemma~\ref{rc:step} at this same prescribed mesh
$\eta$. Impose its passage-and-mass event determined at that crossing,
although only the mass contributions inside $\mathcal D$ will be
counted. The auxiliary later suffix used to prove the lemma is not
retained. Given previous successes, the stopped history is therefore
admissible for the next step. Denote the intersection of these events
by $E_\eta$.
Iterated conditional probabilities give
\[
 \mathbb P(E_\eta)\ge q_{\mathrm{init}}q^N
           \ge c\exp(-Cm^2).
\]
On $E_\eta$, the disjoint middle strips inside the disk yield
normalized sausage area at the chosen mesh $\eta$ of at least
\begin{equation}\label{rc:layout-sum}
 c m^2\ell^d\ge c m^{2-d}=c m^p.
\end{equation}
Disjointness concerns the regions in which area is counted, so it
continues to hold if a sausage also meets another piece of the curve.

On this $\eta$-dependent success event, the sausage comparison in
\eqref{rc:compact-sausage-transfer} maps the area counted in
\eqref{rc:layout-sum} into the physical $v$-sausage. Conformal area
change multiplies normalized area by at least $c r^d$, so
\eqref{rc:dense-mass} follows. The choice of $\delta_m$ above depends
only on $m$ and the fixed constants, uniformly in the stopped state.

It remains to check the radius cutoff and the change of law.  Through
completion, the prescribed trace and its filling remain in a narrow
neighborhood of the connector, array disk and short terminal extension.
More precisely, choose a closed obstacle consisting of $\mathcal D$,
a thin neighborhood of the simple connector, and a thin neighborhood of
the terminal stub from the disk exit through the completion front.
The stub has length at most $2\ell$. The connector and stub attach to
disjoint small arcs of the disk, and their neighborhoods are otherwise
disjoint. This obstacle has
connected complement in $\mathbb H$, so every filled prefix hull of a
trace contained in it is contained in it as well.

The point $w$ is separated from this set and the real boundary by a
fixed positive distance.  The connector is below $w$ until its final
vertical portion, whose horizontal distance from $w$ is bounded below
by a fixed multiple of $s_0$; the array and terminal extension have the
same separation.  Thus $w$ can move vertically to height $2$, then
horizontally to either $x=3$ or $x=-3$, and then down to a small real
interval about that coordinate, with uniformly positive clearance from
the prescribed trace except at the real endpoint.  Small adjustments
of these polygonal routes give fixed-width neighborhoods.  They also
connect $w$ and a fixed ball around it to infinity, so these sets are
in the surviving domain throughout the experiment.  The two chosen
real intervals remain unvisited and unswallowed; their Loewner
coordinates stay on their respective sides of the driving point by
continuity.

Brownian tracking along the two routes therefore gives
a uniform lower bound for the two boundary-arc harmonic measures in
every successful intermediate domain.  Hence the
marked-point angle sine stays at least $s_1>0$.  Its conformal radius
stays above $c_1>0$ in parameter coordinates.  Conformal covariance
therefore keeps its physical radius above $c_2r$.  Choose $C_1>1/c_2$
with strict slack.

Let $T$ be the first crossing of the completion front. It is a stopping
time and is finite on $E_\eta$. On $E_\eta$, the one-point density at
this stop,
$M=R^{-p}\sin^{a_0}\theta$, $a_0=8/\kappa-1$, satisfies
\[
 \frac{M_T}{M_0}
 =\left(\frac{R_T}{r}\right)^{-p}
     \left(\frac{\sin\theta_T}{\sin\theta_0}\right)^{a_0}
 \ge s_1^{a_0}.
\]
Stop the density at the radius cutoff and use the finite stopped
change of law, first also stopping at a deterministic capacity horizon
and then increasing that horizon.  Since success occurs in finite
time and before the radius cutoff, it follows that its tilted
probability is at least $s_1^{a_0}$ times its ordinary probability.
This proves \eqref{rc:dense-probability} and
\eqref{rc:dense-mass}.  All tests concern the already generated
sausage at finite stops; no limiting marked-point visit is used.
\end{proof}
